\chapter{Asking the right question}\label{ch:questions}

The ask-back loop is what makes Near different, and we never cut it. After each report, Near considers
every question it could ask, predicts how each possible answer would change the position, and asks the one
expected to help most.

\section{Candidate questions}

\begin{table}[htbp]\centering\small
\caption{Question types, the answers modelled, and their effort (load) for the person in the field.}\label{tab:qtypes}
\begin{tabularx}{\linewidth}{lXlc}
\toprule
Kind & Example & Answers & Load\\
\midrule
visible\_type & ``Can you see a railway?'' & yes / no & 1\\
visible\_feature & ``Can you see the Sainte-Chapelle?'' & yes / no & 1\\
which\_see & ``Which of these can you see: the chapel, a tower, the Tour des Salves?'' & any subset & 1.5\\
setting3 & ``Inside, in a street, or in the open?'' & 3 & 1.5\\
hear\_type & ``Have you heard a train in the last few minutes?'' & yes / no & 1\\
ground & ``Under trees? Tarmac, dirt, or no path?'' & 2--3 & 1--1.5\\
left\_right & ``Facing the keep, is the church on your left or right?'' & left / right & 2\\
sun & ``Where is the sun? Give me a clock direction.'' & 12 & 2\\
compass & ``If you have a compass: bearing to the keep? Degrees or mils.'' & 16 sectors or none & 2\\
feature\_dir & ``Where is the Sainte-Chapelle from you? Clock direction.'' & 12 or can't see & 2\\
hear\_dir & ``Where do you hear road traffic from?'' & 12 or none & 2\\
dist\_band & ``Roughly how far is the church, in metres?'' & 5 bands & 2.5\\
describe & ``What landmarks can you see? Clock directions.'' & open & 3\\
\bottomrule
\end{tabularx}
\end{table}

Questions are rebuilt every turn from the current likely spots A/B/C and the real named features around
them. The ``which do you see'' question uses the features whose visibility differs most between the spots.

\section{Choosing one}

For a question $q$ with answers $a$, the answer model gives $P(a \mid p, \theta)$ for every cell and
heading, with realistic noise (a person misjudging left and right, a compass answer in 22.5$^\circ$
sectors, a 40 \% chance of having no compass at hand). With the current posterior $P(p,\theta)$:
\begin{align}
  P(a) &= \sum_{p,\theta} P(a \mid p,\theta)\,P(p,\theta), \\
  P(p,\theta \mid a) &\propto P(a \mid p,\theta)\,P(p,\theta), \\
  \Delta H(q) &= H\big[P(p)\big] - \sum_a P(a)\, H\big[P(p \mid a)\big] \quad\text{(bits)}, \\
  \Delta S(q) &= S\big[P(p)\big] - \sum_a P(a)\, S\big[P(p \mid a)\big] \quad\text{(metres)},
\end{align}
where $S$ is the spread of the position (expected distance to the best guess). Bits alone favour
splitting far-off areas the coordinator doesn't care about; metres alone ignore ambiguity. The score
blends both, then divides by the effort:
\begin{equation}
  \text{score}(q) \;=\; \frac{\text{blend}\big(\Delta H(q), \Delta S(q)\big)}{1 + \lambda\,(\text{load}(q) - 1)} .
\end{equation}

\paragraph{Stopping.} Near stops when the radius is under 75 m, when no question gains at least 0.05 bits,
or when the budget is spent: up to 6 location questions (5 when the speaker sounds stressed, 4 when highly
stressed).

\paragraph{A goal for every question.} The dashboard and the phone show why each question is asked: ``Goal:
tell area A (45 \%) from area B (30 \%) · yes $\rightarrow$ A, no $\rightarrow$ B · error $\sim$600 m $\rightarrow$ $\sim$150 m'',
or ``Goal: a compass bearing removes the unknown facing direction''.

\nanofig{questions}{Choosing a question: each candidate is scored by how much it is expected to shrink the position, given
every answer it could get.}
{Decision-tree infographic, 16:9, plain white background, flat vector. At the top a small map with two
blobs labelled ``A 45\%'' and ``B 30\%''. Below, three candidate question cards side by side: ``Can you see a railway?''
with two branches ``yes'' and ``no'' each showing a tiny map where the blobs change little, score ``0.2 bits'';
``Is the church left or right of the keep?'' with branches ``left'' (only A remains) and ``right'' (only B remains),
score ``0.95 bits, 600 m → 150 m'' highlighted with a cyan glow and a check mark; ``How far is the church?'' with five
small branches, score ``0.4 bits, effort high''. A small label ``asked'' over the winning card. Exactly these labels. Colour palette (use only these): navy \#0D1110, cyan \#1AA7C4, amber \#E09A2B, red \#E5484D, green \#3FA877, grey \#8B9A94, plus white for text on dark shapes.}

\section{Adapting to stress}

\code{stress.py} measures each push-to-talk clip within about 0.3 s of release. It looks at pitch raised
above the speaker's own floor, pitch range (its own numpy pitch tracker, $\pm$2 Hz on test tones),
brightness (1--5 kHz against 50 Hz--1 kHz), articulation rate, choppy speech, and words such as ``hurry'',
``he's dying'', \emph{vite}, repeats and exclamations. The result is low, elevated or high. Shouting over
noise raises pitch and brightness too, so those cues count less when the background is loud. Stress only
changes \emph{how} Near asks. It never touches medical priority.

\begin{table}[htbp]\centering\small
\caption{Question policy by stress level.}\label{tab:stress}
\begin{tabular}{llll}
\toprule
Speaker & Hardest question allowed & Weight $\lambda$ on effort & Wording\\
\midrule
low / unknown & describe everything (3) & 0.10--0.15 & normal\\
elevated & distances (2.5) & 0.35 & normal\\
high & one clock direction or left/right (2) & 0.60 & shortest (``Sun at what o'clock?''), slower read-back\\
\bottomrule
\end{tabular}
\end{table}

\section{Medical questions}

Health questions are asked in MARCH order and don't count against the location budget. Life threats come
first: is the bleeding stopped, is the casualty awake and answering, are they breathing normally, and where
is the wound when an injury has no body part. Then location and health questions take turns: how it
happened, pale, cold or sweaty, how many minutes ago. When the speaker is the casualty, Near says ``you''
(``Has your bleeding stopped?'').

\section{Results}

\begin{table}[htbp]\centering\small
\caption{Question objective, 40 simulated street casualties with 3D line of sight, same spots and draws in every row
(\code{scripts/eval\_questions.py}).}\label{tab:qobj}
\begin{tabular}{lrrr}
\toprule
& Random street spots & \multicolumn{2}{c}{Spots that see a landmark}\\
\cmidrule(lr){3-4}
Objective & median error & median error & within 250 m\\
\midrule
bits only (before) & 735 m & 359 m & 42 \%\\
metres only & 666 m & 344 m & 45 \%\\
\textbf{blend (current)} & 669 m & \textbf{201 m} & \textbf{52 \%}\\
blend, elevated stress & 666 m & 276 m & 50 \%\\
blend, high stress (3 easy questions) & n/a & 617 m & 35 \%\\
\bottomrule
\end{tabular}
\end{table}

Easier questions cost accuracy. Under high stress, the best setting we tested keeps 3 quick questions (617 m);
cutting to 2 gave 791 m. That trade between speed and precision is for the medical mentors to decide.
